\chapter {Metodologia}

Este capitulo explica a metodologia adotada neste trabalho, descrevendo suas etapas, ferramentas utilizadas e os perfis envolvidos na resolução da problemática central. O trabalho segue o formato de estudo de caso, com características de pesquisa aplicada mista.

\section {Síntese do Referencial Teórico}

A fundamentação teórica presente neste trabalho foi feita com objetivo de trazer o conhecimento necessário para completo entendimento dos tópicos envolvidos na solução do problema de desempenho e adequação da ferramenta de Textos Participativos do Brasil Participativo. Os tópicos envolvidos na execução deste trabalho incluem desde técnicas organizacionais de desenvolvimento de \textit{software} até a abordagens de desenvolvimento sob o rigor de requisitos de desempenho mais restritos, especialmente com o uso do \textit{framework Ruby on Rails}.

\section {Ferramentas}

Para a execução desse trabalho foi adotado o uso de diversas ferramentas:

\subsection {Visual Studio Code}

O Visual Studio Code (VSCode) é um \textit{software} editor de textos de código aberto amplamente utilizado para codificação de sistemas e aplicações. No VSCode é possível trabalhar com o repositório de código da aplicação de forma fácil e integrada, com capacidades de \textit{syntax highlighting}, \textit{auto-formatting} e visualizações do gerenciador de versões Git. Esse editor ainda conta com uma grande gama de extensões da comunidade para facilitar e melhorar a qualidade do desenvolvimento. Uma das extensões utilizadas foi o RubyLSP, que permite a verificação do código em tempo real, capturando erros de sintaxe, além de fornecer um mecanismo de crítica do código considerando as regras de \textit{lint} definidas no projeto.

\subsection {Git}

O Git é uma ferramenta gratuita de código aberto de gestão de versionamento distribuído de código, que pode ser utilizado desde pequenas até grandes aplicações. O Git é amplamente utilizado para gerenciamento de projetos modernos de código.

\subsection {Yarn}

O Yarn é um gerenciador de pacotes de código aberto para gerenciar a distribuição de dependência em projetos JavaScript. Seu objetivo é ajudar no processo de instalação, atualização, configuração e remoção de dependências de pacotes.

\subsection {RuboCop}

O RuboCop é um analisador de estilo de código e formatador \textit{Ruby} que se baseia no guia de estilo de código \textit{Ruby Style Guide}, que é guiado pela comunidade de desenvolvedores. O uso do RuboCop juntamente do RubyLSP acelerou o desenvolvimento e a legibilidade do código, trazendo em tempo real sugestões de refatoração de acordo com os padrões da comunidade.

\subsection {Nokogiri}

O Nokogiri é uma \textit{gem Ruby} usada para manipular documentos HTML e XML. A biblioteca conta com uma conjunto grande de funções que auxiliam o processamento de HTML, permitindo com que o desenvolvedor não precise se preocupar com problemas específicos de representação de XML e HTML no \textit{Ruby}.

\subsection {Jodit}

O Jodit é um editor de textos para navegadores. Este editor segue o formato \textit{What You See is What You Get} (WYSIWYG), que basicamente permite com que o usuário insira o texto com a formatação desejada e com a confiança de que essa formatação será persistida e exibida posteriormente.

\subsection {LibreOffice Writer}

O LibreOffice Writer é uma ferramenta de edição de textos de código aberto similar ao Microsoft Word. O LibreOffice Writer permite a escritra de textos, exportação em formato PDF e afins. O LibreOffice Writer é utilizado para trabalhar com os arquivos do formato \textit{docx}, que é utilizado para a primeira confecção dos Textos Participativos pelos usuários.

\subsection {Webpack e Webpacker}

O Webpack é uma ferramenta de empacotamento escrita em JavaScript usada para realizar a transformação de ativos servidos pela aplicação para um formato compreendido pelos navegadores web. Os ativos da aplicação podem incluir desde imagens, arquivos de estilo até código JavaScript. O Webpacker, por sua vez, é uma extensão disponível para o \textit{Ruby on Rails} que realiza a interface de comunicação entre a aplicação, que é escrita em Ruby, com o Webpack.

\subsection {Notion}

Em geral, o Notion pode ser definido como uma aplicação \textit{web} utilizada para anotações. O Notion conta com a possibilidade de se trabalhar em conjunto em espaços de trabalho colaborativos. O Notion pode servir como um repositório de ideias e anotações na nuvem, facilitando a recuperação e organização.

\subsection {ApacheBench}

O ApacheBench é uma ferramenta de linha de comando utilizado para realizar testes de carga em aplicações \textit{web}. Essencialmente ela serve para mostrar quantas requisições por segundo a aplicação é capaz de entregar. Sua simplicidade de uso é um de seus pontos fortes.

\subsection {Draw.io}

O Draw.io é uma ferramenta de criação de diagamas online gratuita, permitindo a criação de diagramas clássicos como UML, ER e afins.

\section {Abordagem}

A abordagem deste trabalho pode ser dividida em duas vertentes: a abordagem organizacional, e a abordagem técnica. A abordagem organizacional diz respeito à forma como o trabalho foi conduzido durante o tempo, e a divisão das atividades. A abordagem técnica demonstra como os problemas de \textit{software} foram tratados para se chegar no resultado esperado.

\subsection {Abordagem Organizacional}

A fim de alcançar os objetivos destacados foi adotada a metodologia de desenvolvimento de \textit{software} ágil (\textit{agile}). A rotina de execução de trabalho foi dividida em \textit{sprints}, onde em cada uma delas um conjunto de funcionalidades e atividades serão realizadas. No fim de cada \textit{sprint}, o que foi feito é validado com o time de produto e se necessário com os \textit{stakeholders} do projeto.

\subsection {Abordagem Técnica}

A abordagem técnica será o desenvolvimento orientado à busca por desempenho. Nesse contexto, partes do sistema que devem receber ajustes serão previamente analisadas em termos de desempenho através da checagem das requisições usando da própria ferramenta de \textit{logs} do \textit{framework}. Uma boa ferramenta de logs fornecerá ao desenvolvedor detalhes de tempo gasto por cada componente chave do sistema. Nesse momento é importante que as mesmas condições do ambiente produtivo sejam replicadas no ambiente local de desenvolvimento sempre que possível.

Caso os ajuste em uma determinada parte da funcionalidade envolva a melhoria de desempenho, após a realização das modificações, a mesma análise é conduzida para validar a qualidade dos ajustes e seguir adiante.

\section {Compilação de Resultados}

Como artefato final deste trabalho, será feita a compilação dos resultados obtidos com os ajustes realizados na plataforma do Brasil Participativo sobre o uso da ferramenta de Textos Participativos. Para demonstrar de forma quantitativa os resultados alcançados, será feita uma tabela comparando os dados de desempenho da aplicação em testes de simulação realizados com a ferramenta ApacheBench.